% -- math shorthands ----------------------------------------------------------

\newcommand{\F}{\mathbb{F}}
\newcommand{\Z}{\mathbb{Z}}
\newcommand{\N}{\mathbb{N}}
\newcommand{\Q}{\mathbb{Q}}
\newcommand{\R}{\mathbb{R}}
\newcommand{\C}{\mathbb{C}}

\newcommand{\set}[1]{\left\{#1\right\}}
\newcommand{\abs}[1]{\left|#1\right|}
\newcommand{\norm}[1]{\left\lVert #1\right\rVert}
\newcommand{\inner}[2]{\left\langle #1,#2\right\rangle}
\newcommand{\gen}[1]{\left\langle #1\right\rangle}
\newcommand{\ideal}[1]{\left\langle #1\right\rangle}
\newcommand{\vect}[1]{\mathbf{#1}}
\newcommand{\mat}[1]{\mathbf{#1}}

% -- coding theory shorthands -------------------------------------------------

\newcommand{\code}{\mathcal{C}}
\newcommand{\wt}{\operatorname{wt}}
\newcommand{\rank}{\operatorname{rank}}
\newcommand{\Span}{\operatorname{span}}
\newcommand{\supp}{\operatorname{supp}}
\newcommand{\poly}{\operatorname{poly}}
\newcommand{\negl}{\operatorname{negl}}
\newcommand{\ol}[1]{\overline{#1}}

\newcommand{\nkcode}[3]{[#1,#2]_{#3}}
\newcommand{\nkdcode}[4]{[#1,#2,#3]_{#4}}

% -- presentation teaching lens -----------------------------------------------

\newcommand{\hardproblemlens}{%
  Code-based hardness lives in the gap between linear algebra and combinatorial search.%
}
\newcommand{\HardProblemLens}[1]{%
  \gdef\hardproblemlens{#1}%
}

% -- colors -------------------------------------------------------------------

\definecolor{accent}{HTML}{1F4E79}
\definecolor{softaccent}{HTML}{EAF2F8}
\definecolor{softgray}{HTML}{F6F7F9}
\definecolor{successdraw}{HTML}{27AE60}
\definecolor{successfill}{HTML}{EAFAF1}
\definecolor{errordraw}{HTML}{C0392B}
\definecolor{errorfill}{HTML}{FDEDEC}

\definecolor{Ink}{HTML}{101A24}
\definecolor{Paper}{HTML}{F8FBFC}
\definecolor{Grid}{HTML}{DCEAEC}
\definecolor{BlueCode}{HTML}{3155A4}
\definecolor{AmberCode}{HTML}{B87511}
\definecolor{ErrorRed}{HTML}{D33F49}
\definecolor{GoodGreen}{HTML}{3C8D4F}
\definecolor{VioletCode}{HTML}{6E5AA8}

% -- tikz styles --------------------------------------------------------------

\tikzset{
  every picture/.style={
    font=\sffamily,
    line cap=round,
    line join=round
  },
  title text/.style={
    font=\sffamily\bfseries\large,
    text=TealDarkCyan
  },
  caption text/.style={
    font=\sffamily\small,
    text=Ink!72
  },
  tiny label/.style={
    font=\sffamily\scriptsize,
    text=Ink!72
  },
  bit/.style={
    draw=TealDarkCyan!55,
    fill=Paper,
    minimum width=0.46cm,
    minimum height=0.46cm,
    inner sep=0pt,
    font=\sffamily\scriptsize\bfseries
  },
  bit teal/.style={
    bit,
    fill=AccentCyan!13,
    draw=AccentCyan!75,
    text=TealDarkCyan
  },
  bit blue/.style={
    bit,
    fill=BlueCode!12,
    draw=BlueCode!72,
    text=BlueCode!85!black
  },
  bit amber/.style={
    bit,
    fill=AmberCode!15,
    draw=AmberCode!72,
    text=AmberCode!85!black
  },
  bit error/.style={
    bit,
    fill=ErrorRed!13,
    draw=ErrorRed!80,
    text=ErrorRed!85!black
  },
  panel/.style={
    draw=TealDarkCyan!26,
    fill=Paper,
    rounded corners=3pt,
    line width=0.8pt,
    inner sep=6pt
  },
  process/.style={
    draw=TealDarkCyan!68,
    fill=TealDarkCyan!7,
    rounded corners=3pt,
    line width=0.9pt,
    align=center,
    inner sep=6pt
  },
  process blue/.style={
    process,
    draw=BlueCode!70,
    fill=BlueCode!8,
    text=BlueCode!80!black
  },
  process amber/.style={
    process,
    draw=AmberCode!75,
    fill=AmberCode!10,
    text=AmberCode!80!black
  },
  process red/.style={
    process,
    draw=ErrorRed!75,
    fill=ErrorRed!9,
    text=ErrorRed!85!black
  },
  flow/.style={
    -{Stealth[length=3.0pt]},
    line width=0.95pt,
    draw=TealDarkCyan!80
  },
  flow blue/.style={
    flow,
    draw=BlueCode!80
  },
  flow amber/.style={
    flow,
    draw=AmberCode!82
  },
  flow red/.style={
    flow,
    draw=ErrorRed!80
  }
}
